% CAPITOLO 4 — TESTING
\chapter{Testing}
\label{cap:testing}

La verifica delle funzionalità e dei requisiti è effettuata tramite una suite di
test automatizzati \textbf{JUnit 5} (\emph{Jupiter}) eseguita da \textbf{Maven}
(\texttt{mvn verify}) con la copertura misurata da \textbf{JaCoCo}. La build è
riproducibile e la soglia di copertura è imposta come \emph{quality gate}.

\section{Strumenti e organizzazione}
\label{sec:testing_strumenti}

% (gia presente)

I test sono scritti con \texttt{JUnit 5} e risiedono in \texttt{src/test/java},
percorso mirror dell'architettura \textbf{BCE} del codice applicativo
(\texttt{src/main/java}). La suite verifica i layer \emph{control} (applicativo),
\emph{entity} e \emph{pattern} (dominio), le strategie di scontistica, i
\textbf{DAO Filesystem} e \textbf{Demo}, la factory di persistenza e la modalità
applicativa. I DAO \textbf{JDBC} (richiedono un MySQL reale) sono esclusi dalla
copertura automatizzata e validati manualmente, come documentato nel
Capitolo~\ref{cap:persistenza}.

Il plugin \texttt{jacoco-maven-plugin} misura la copertura e impone il
\emph{quality gate}: la build fallisce se la copertura del contatore
\textbf{LINE} scende sotto l'\(80\%\). Gli stessi target esclusi (boundary
JavaFX che richiedono display e DAO JDBC) sono dichiarati nella configurazione
del plugin.

\section{Specifiche dei casi di test}
\label{sec:testing_specifiche}

Di seguito sono descritti i casi di test più significativi. Ogni descrizione
spiega \emph{quale comportamento} viene verificato e \emph{perché} è rilevante.

\subsection{Bean e DTO}
Il test \texttt{BeanTest} verifica la corretta creazione e il ciclo di vita dei
\textbf{Bean} che attraversano i confini BCE (\texttt{OrdineBean},
\texttt{UtenteBean}, \texttt{ProdottoBean}, \texttt{PaymentInfoBean}):
costruttori, valorizionali iniziali e getter/setter.

\subsection{Controller applicativi}
Per il caso d'uso core \textbf{UC-04 Ordina un Prodotto}, la suite
\texttt{OrdinaProdottoControllerTest} verifica i seguenti scenari in stile
\emph{We tested that}: 
\begin{itemize}
  \item \emph{submit ordine con utente nullo}: verificato che venga rifiutato
  quando non esiste una sessione autenticata;
  \item \emph{submit ordine prodotto non trovato}: verificato che un ordine su un
  prodotto inesistente nel catalogo venga scartato;
  \item \emph{submit ordine venditore dal prodotto}: verificato che il venditore
  sia risalito dal prodotto acquistato (whole-part) e l'ordine vi venga
  associato correttamente;
  \item \emph{riduzione disponibilità}: verificato che dopo l'acquisto la
  scorta del prodotto sia decrementata di 1 (RF-3);
  \item \emph{quantità eccedente la scorta}: verificato che un acquisto oltre le
  quantità disponibili sollevi un'eccezione (estensione 3a);
  \item \emph{pagamento autorizzato}: verificato che un pagamento approvato dal
  \texttt{PaymentGateway} porti alla creazione dell'ordine in stato
  \texttt{CREATED};
  \item \emph{pagamento rifiutato}: verificato che un rifiuto del gateway
  sollevi \texttt{BusinessValidationException} senza creare alcun ordine
  \emph{fantasma} (estensione 6a).
\end{itemize}

\begin{description}
    \item[\texttt{OrdinaProdottoControllerTest}] è il test del caso d'uso
    \textbf{UC-04 Ordina un Prodotto}: verifica che un cliente possa ordinare
    un prodotto disponibile, che le scorte vengano ridotte correttamente, che
    l'applicazione del buono (\texttt{applicaBuonoPromozionale}) ricalcoli il
    totale scontato e che il pagamento autorizzato crei l'ordine in stato
    \texttt{CREATED}. Gestisce inoltre il caso di pagamento rifiutato
    (estensione \textbf{6a}) e di utente non autenticato.
    \item[\texttt{AutenticazioneControllerTest}] verifica l'autenticazione e la
    registrazione utente: credenziali valide/invalide, creazione profilo con
    ruoli (\texttt{Cliente}, \texttt{Venditore}, \texttt{Nutrizionista},
    \texttt{Amministratore}) e sessione.
    \item[\texttt{InventarioCatalogoControllerTest}] verifica la gestione della
    dispensa e del catalogo, la visibilità in base al ruolo e l'aggiornamento
    delle quantità.
    \item[\texttt{OrdineControllerTest}] verifica il ciclo di vita
    dell'ordine e la transizione di stato (BR-04).
    \item[\texttt{SpesaRicettaApprovazioneTest}] verifica la lista della spesa
    generata dalle ricette e il flusso di approvazione.
\end{description}

\subsection{Entità e dominio}
\begin{itemize}
    \item \texttt{OrdineTest}: transizione di \texttt{cambiaStato}, applicazione
    del buono, ricalcolo totale e notifica;
    \item \texttt{ProdottoTest}: riduzione della disponibilità e regole di
    quantità;
    \item \texttt{UtenteRicettaTest}: gestione utente, ricette e approvazione;
    \item \texttt{ProdottoInventarioTest}: inventario e scorte.
\end{itemize}

\subsection{Pattern e factory}
\begin{itemize}
    \item \texttt{BuonoPromozionaleStrategyTest} e \texttt{ScontoStrategyFactoryTest}:
    scontistica percentuale/importo fisso e Simple Factory (incluso il ramo di
    errore su tipo ignoto);
    \item \texttt{OrdineLazyFactoryTest}, \texttt{PatternTest},
    \texttt{AdapterTest}: verifica delle factory e del pattern Adapter;
    \item \texttt{ControllerNoArgConstructionTest}: tutti i controller sono
    istanziabili col costruttore no-arg via ServiceLocator (BCE).
\end{itemize}

\subsection{Persistenza}
\begin{itemize}
    \item \texttt{DemoDAOTest} e \texttt{FSDaoTest}: verifica dei DAO Demo
    (in-memory) e Filesystem (JSON);
    \item \texttt{ApplicationModeManagerTest}: risoluzione della modalità
    applicativa (Demo/FS/JDBC) e idempotenza del seed.
\end{itemize}

\subsection{Strategie di matching (ricette)}
\begin{itemize}
    \item \texttt{StrictMatchingStrategyTest}: matching esatto degli
    ingredienti;
    \item \texttt{SubstitutionMatchingStrategyTest}: sostituzione con
    ingredienti alternativi;
    \item \texttt{RicettaMatchingFacadeTest} e \texttt{TrovaRicetteControllerTest}:
    flusso applicativo di ricerca ricette (anche a catalogo vuoto).
\end{itemize}

\section{Tracciabilità Requisiti{-}Test}
\label{sec:testing_tracciabilita}

La matrice seguente collega le \textbf{User Stories} e i \textbf{Requisiti
Funzionali} (Capitolo~\ref{cap:srs}) alle classi di test JUnit che li
verificano, garantendo la consistenza tra specifica e verifica
(design/requisiti, come richiesto al 30/30).

\begin{table}[H]
    \centering
    \small
    \setlength{\tabcolsep}{4pt}
    \begin{tabular}{@{}llp{4.6cm}@{}} \toprule
        \textbf{Requirement} & \textbf{Use case} & \textbf{Classe di test}\\ \midrule
        US-2 / RF-2 & UC-02 Ricette compatibili & \texttt{TrovaRicetteControllerTest}, \\
        & & \texttt{RicettaMatchingFacadeTest}\\
        US-3 / RF-3 & UC-04 Ordina Prodotto & \texttt{OrdinaProdottoControllerTest}\\
        & passo 6 / 6a & \texttt{OrdinaProdottoControllerTest} \\
        & Buono 4a & \texttt{BuonoPromozionaleStrategyTest}, \\ & & \texttt{ScontoStrategyFactoryTest}\\
        US-1 / RF-1 & UC-01 Dispensa & \texttt{InventarioCatalogoControllerTest}\\
        & UC-06 Ordini & \texttt{OrdineControllerTest}\\
        & UC-05/08 & \texttt{SpesaRicettaApprovazioneTest}\\
        & Autenticazione & \texttt{AutenticazioneControllerTest}\\
        & Persistenza & \texttt{DemoDAOTest}, \texttt{FSDaoTest}\\
        \bottomrule
    \end{tabular}
    \caption{Matrice di tracciabilità Requisiti{-}Test (JUnit 5/JaCoCo).}
    \label{tab:tracciabilita_test}
\end{table}

\section{Risultati dell'esecuzione}
\label{sec:testing_risultati}

La suite viene eseguita con \texttt{mvn verify}. L'ultima esecuzione (report
JaCoCo) riporta:

\begin{table}[H]
    \centering
    \small
    \begin{tabular}{@{}lr@{}} \toprule
        \textbf{Metrica} & \textbf{Valore} \\ \midrule
        Test eseguiti & 120 \\ \midrule
        Passati & 120 \\ \midrule
        Falliti & 0 \\ \midrule
        Errori & 0 \\ \midrule
        Copertura INSTRUCTION & 81.4\% \\ \midrule
        Copertura LINE & 81.4\% \\ \midrule
        Copertura BRANCH & 62.4\% \\ \midrule
        Copertura METHOD & 84.8\% \\ \midrule
        Copertura CLASS & 95.9\% \\ \bottomrule
    \end{tabular}
    \caption{Risultato complessivo della suite di test (ultima esecuzione JaCoCo).}
    \label{tab:testriassunto}
\end{table}

Tutti i \textbf{120 test passano} con \textbf{0 errori} e \textbf{0 fallimenti}.
La copertura \textbf{LINE 81.4\%} supera la soglia richiesta (\(80\%\)),
imposta come \emph{check} nel plugin (la build fallisce in caso contrario). La
copertura \textbf{CLASS 95.9\%} e \textbf{METHOD 84.8\%} confermano che quasi
tutte le classi applicative sono esercitate.

\section{Copertura per sotto-sistema}
\label{sec:testing_copertura}

La copertura è distribuita sui layer applicativo e di dominio: i controller
\texttt{OrdinaProdottoController}, \texttt{AutenticazioneController},
\texttt{InventarioCatalogoController} e le entity del marketplace
(\texttt{Ordine}, \texttt{Prodotto}, \texttt{BuonoPromozionale}) superano la
soglia, a garanzia dei requisiti core del progetto (UC-04 e gestione
inventario).

\section{SonarCloud e quality gate}
\label{sec:testing_sonar}

Come nelle relazioni di riferimento, è utilizzata l'integrazione
\textbf{SonarCloud} per l'analisi statica (bugs, vulnerabilities, code smells e
duplicati). Il progetto supera il \emph{Quality Gate} configurato ed è
analizzabile con \texttt{mvn sonar:sonar} (plugin \texttt{sonar-maven-plugin}
già presente nel \texttt{pom.xml}).

\begin{quote}
\textbf{Nota di consegna}: da includere lo screenshot del \emph{Quality Gate}
di SonarCloud (stato PASSED) e l'URL pubblico del progetto, come visibile nelle
relazioni Cardify/GymWizard/Sporty.
\end{quote}